\title{Controlling Text-to-Image Diffusion by Orthogonal Finetuning}

\begin{abstract}
Large-scale text-to-image diffusion models have demonstrated remarkable prowess in synthesizing photorealistic images from textual descriptions. However, devising effective strategies to steer or constrain these models for various downstream applications remains a significant unsolved problem. In this work, we propose Orthogonal Finetuning~(OFT), a systematic approach for adapting text-to-image diffusion models to specific downstream objectives. Distinct from prior techniques, OFT is theoretically guaranteed to maintain hyperspherical energy, a metric reflecting the pairwise interactions between neurons on a unit hypersphere. Preserving this quantity proves essential for retaining the semantic image generation capabilities inherent in text-to-image diffusion models. To further enhance the stability of the finetuning process, we introduce Constrained Orthogonal Finetuning (COFT), which incorporates an additional constraint on the radius of the hypersphere. Our investigation focuses on two core finetuning scenarios in text-to-image tasks: subject-driven generation, which entails synthesizing subject-specific images given a handful of samples and a text prompt, and controllable generation, where models must incorporate additional structured control cues. Through empirical evaluation, we demonstrate that our OFT approach consistently surpasses existing alternatives in both image generation quality and speed of convergence.
\end{abstract}

\section{Introduction}

State-of-the-art text-to-image diffusion models~\cite{saharia2022photorealistic,ramesh2022hierarchical,rombach2022high} have recently achieved unprecedented results in generating high-quality images based on textual instructions. Nonetheless, relying solely on textual guidance often leaves room for ambiguity and lacks the ability to impose precise, fine-grained control over the resulting images. To overcome these limitations, our study focuses on two primary categories of text-to-image generation tasks:

\begin{itemize}[leftmargin=*,nosep]

    \item \textbf{Subject-driven generation}~\cite{ruiz2023dreambooth}: The objective is to create new images depicting a specific subject, given a few example images and accompanying text prompts that specify context.
    \item \textbf{Controllable generation}~\cite{zhang2023adding,mou2023t2i}: This involves utilizing auxiliary control inputs (such as canny edge maps or segmentation masks) in combination with text prompts, directing the model to generate images that faithfully follow these controls.
\end{itemize}

The core challenge in both tasks involves determining how to optimize the finetuning process for text-to-image diffusion models so as to retain the original generative capabilities established during pretraining. Key requirements for a robust finetuning approach are: (1) \emph{training efficiency}, characterized by a reduction in both the number of trainable parameters and the training epochs needed, and (2) \emph{generalizability preservation}, ensuring the model maintains its ability to generate high-quality and diverse outputs. Typically, this is accomplished either by gradually modifying neuron weights using a small learning rate (\emph{e.g.}, \cite{ruiz2023dreambooth}) or by integrating a minor component featuring re-parameterized neuron weights (\emph{e.g.}, \cite{hulora2022,zhang2023adding}). However, these straightforward strategies do not offer guarantees regarding the retention of pretraining generative performance. Additionally, there is a notable lack of systematic principles for developing effective finetuning methods or for identifying optimal hyperparameters like the training duration. A significant obstacle arises from the absence of reliable metrics to assess the extent to which the pretrained generative abilities are preserved. Many current finetuning approaches operate under the implicit assumption that reducing the Euclidean distance between the finetuned and original models leads to better retention of pretrained performance. For this reason, finetuning is often conducted with an exceedingly low learning rate. Nevertheless, while this premise might be valid in certain scenarios, we contend that relying solely on Euclidean proximity does not adequately reflect the preservation of semantic content. Therefore, adopting a more structurally informative metric to measure the divergence between finetuned and pretrained models could substantially enhance the ability to safeguard pretraining performance and promote stable finetuning outcomes.

Drawing on empirical findings that hyperspherical similarity effectively captures semantic structure~\cite{liu2017deep,liu2018decoupled,chen2020angular}, our approach leverages hyperspherical energy~\cite{liu2018learning} to quantify the inter-neuron relational patterns. Hyperspherical energy is calculated by summing the hyperspherical similarity (\emph{e.g.}, cosine similarity) across all neuron pairs within a given layer, thereby measuring how uniformly neurons are distributed on the unit hypersphere~\cite{liu2021learning}. We posit that optimal finetuning should preserve the hyperspherical energy, resulting in minimal discrepancy between the finetuned and original pretrained models. Although a straightforward strategy might introduce a regularization term to keep the hyperspherical energy constant during finetuning, such an approach does not guarantee minimization of the energy difference. Instead, we capitalize on a key invariance property of hyperspherical energy: applying a uniform orthogonal transformation to all neurons leaves the pairwise hyperspherical similarity intact. Building on this property, we introduce Orthogonal Finetuning~(OFT), a method that modifies large text-to-image diffusion models for specific tasks without altering the original hyperspherical energy. The technique centers on learning a single orthogonal transformation per layer that preserves the angular relationships among neurons, effectively rotating the neuron basis in a coordinated fashion. In essence, OFT can be interpreted as redefining the coordinate axes for the neurons within each layer. Considering also that a smaller Euclidean deviation between the finetuned and pretrained models correlates with better retention of pretrained capabilities, we develop an extended method—Constrained Orthogonal Finetuning~(COFT)—which restricts the parameters of the finetuned model to remain within a fixed-radius hypersphere centered at the pretrained neuron positions.

The underlying reasoning for employing orthogonal transformations in neuron finetuning is partly inspired by the properties of the 2D Fourier transform, which enables the decomposition of an image into its magnitude and phase spectra. Notably, the phase spectrum, representing the angular relationship between the input and the basis functions, is responsible for encoding most of the semantic information. For instance, even when the phase spectrum from an original image is combined with a randomly chosen magnitude spectrum, the reconstructed image largely retains its original semantics. This insight implies that modifying neuron orientations plays a central role in altering the semantic content of generated images—a primary objective in both subject-driven and controllable image synthesis. Nevertheless, if neuron directions are adjusted with excessive freedom, it risks undermining the pretrained model's generative capabilities. To mitigate this, we advocate for maintaining the angles between all pairs of neurons, grounded in the belief that these angles are pivotal for neural network knowledge representation. Building on this idea, we propose to learn an orthogonal transformation shared across all neurons in a given layer, ensuring that the hyperspherical energy remains constant.

Further motivation comes from the work on orthogonal over-parameterized training~\cite{liu2021orthogonal}, where classification models are trained from random initialization by means of orthogonal transformations. Randomly initialized neural networks are known to possess, in expectation, low hyperspherical energy, and the objective in~\cite{liu2021orthogonal} is to sustain low energy during training, as reduced hyperspherical energy is linked with improved classification generalization~\cite{liu2018learning,lin2020regularizing}. The study in~\cite{liu2021orthogonal} demonstrates that orthogonal transformations can provide sufficient expressivity for training neural networks capable of generalizing on classification tasks. By contrast, our approach targets the finetuning of text-to-image diffusion models to enhance both controllability and the fidelity of downstream generative performance. The distinction between OFT and the approach of~\cite{liu2021orthogonal} lies in two main areas. First, whereas~\cite{liu2021orthogonal} pursues the minimization of hyperspherical energy, our OFT strategy seeks to maintain the same hyperspherical energy as the pretrained network, thereby safeguarding the pretrained structural properties during finetuning. Reducing the hyperspherical energy in diffusion model finetuning risks erasing intrinsic semantic configurations. Second, OFT explicitly restricts deviations from the pretrained parameters, leading to a constrained formulation; in contrast,~\cite{liu2021orthogonal} adopts no such regularization. Successful finetuning necessitates balancing flexibility and stability, and we contend that our OFT framework attains this balance effectively. Our main contributions can be summarized as follows:

\begin{itemize}[leftmargin=*,nosep]

    \item We introduce Orthogonal Finetuning, a new approach for finetuning text-to-image diffusion models with enhanced controllability. To further promote robust training, we also design a constrained version that restricts the angular displacement from the original pretrained model.
    \item Distinct from prior finetuning strategies, OFT guarantees the conservation of hyperspherical energy during model updates, a property we empirically identify as crucial for maintaining the generative semantics inherent in the pretrained model.
    \item We implement OFT on two tasks: subject-driven generation and controllable generation. Extensive experiments reveal substantial gains over existing methods in generation fidelity, convergence rate, and training stability. Furthermore, OFT is more sample-efficient, achieving successful finetuning with considerably fewer images and epochs.
    \item For the task of controllable generation, we present a novel metric, termed control consistency, which quantitatively measures controllability. This metric involves estimating the control signal from the generated output and comparing it to the initial control signal, offering empirical evidence of OFT’s superior controllability.
\end{itemize}

\section{Related Work}

\textbf{Text-to-image diffusion models}. Recent years have witnessed remarkable advancements~\cite{saharia2022photorealistic,ramesh2022hierarchical,rombach2022high,gu2022vector,nichol2022glide} in the field of text-to-image generation, a progress primarily driven by breakthroughs in diffusion-based generative frameworks~\cite{ho2020denoising,song2021score,song2021denoising,dhariwal2021diffusion} alongside innovations in vision-language representation learning~\cite{su2020vlbert,lu2019vilbert,radford2021learning,wang2021simvlm,li2022blip,li2023blip,alayrac2022flamingo,schuhmann2022laion}. Architectures like GLIDE~\cite{nichol2022glide} and Imagen~\cite{saharia2022photorealistic} utilize pixel-space diffusion modeling; GLIDE employs joint training of its text encoder and diffusion prior with paired data, while Imagen leverages a fixed, pretrained text encoder. In contrast, Stable Diffusion~\cite{rombach2022high} and DALL-E2~\cite{ramesh2022hierarchical} operate in the latent space: Stable Diffusion relies on VQ-GAN~\cite{esser2021taming} to define its visual codebook, whereas DALL-E2 uses CLIP~\cite{radford2021learning} to build a multimodal embedding space shared between text and images. Beyond diffusion-based approaches, research has also explored generative adversarial networks~\cite{reed2016generative,zhang2017stackgan,xu2018attngan,li2019controllable} and autoregressive architectures~\cite{ramesh2021zero,ding2021cogview,wu2022nuwa,yuscaling} for text-to-image generation tasks. Notably, OFT provides a general-purpose finetuning solution compatible with any text-to-image diffusion model.

\textbf{Subject-driven generation}. Techniques to avoid altering the subject include incorporating explicit user-specified masks~\cite{avrahami2022blended,nichol2022glide} as conditional inputs. Other works employ inversion mechanisms~\cite{choi2021ilvr,dhariwal2021diffusion,ramesh2022hierarchical,gal2022image} to enable contextual changes without modifying the central subject. The method of~\cite{hertz2022prompt} achieves both local and global modifications without the need for explicit masks. However, these strategies often struggle to preserve detailed identity traits. Pivotal Tuning~\cite{roich2022pivotal} introduces generator finetuning in the vicinity of an inverted latent vector, combined with regularization to retain subject identity. Similarly,~\cite{nitzan2022mystyle} develops a personalized generative prior from a subject’s facial images. Although~\cite{casanova2021instance} produces multiple instance variations, it sometimes compromises specificity. DreamBooth~\cite{ruiz2023dreambooth}, using a user-defined token and limited subject images, adapts diffusion models by minimizing a reconstruction loss as well as a class-conditional prior preservation loss. Our proposed OFT builds on the DreamBooth paradigm but replaces basic small-step finetuning with a regime based on orthogonal transformations.

\textbf{Controllable generation}. Image-to-image translation is fundamentally an instance of controllable generation. Historically, this problem has primarily been tackled using conditional generative adversarial networks~\cite{isola2017image,zhu2017toward,wang2018high,park2019semantic,choi2018stargan}. More recently, diffusion models have been employed to address image-to-image translation tasks~\cite{saharia2022palette,voynov2022sketch,wang2022pretraining}. Notably, ControlNet~\cite{zhang2023adding} enhances the controllability of pretrained diffusion models by fine-tuning them to accommodate extra control signals, achieving state-of-the-art performance in guided image synthesis. In parallel, T2I-Adapter~\cite{mou2023t2i} similarly improves controllability by fine-tuning pretrained diffusion models. Adopting the experimental protocols from \cite{zhang2023adding,mou2023t2i}, we utilize OFT for fine-tuning pretrained diffusion models, which delivers superior controllability even with less training data and a reduced number of parameters for adaptation. Importantly, OFT incurs no additional computational costs at inference time.

\textbf{Model finetuning}. Adapting large pretrained models to specific downstream tasks via finetuning has become increasingly prevalent~\cite{devlin2018bert,brown2020language,he2022masked}. Adaptation strategies such as those in~\cite{hulora2022,houlsby2019parameter,pfeiffer2020adapterfusion}—widely explored in NLP—represent one major approach. LoRA~\cite{hulora2022}, a method closely related to OFT, leverages a low-rank parameterization for additive weight updates in the fine-tuning stage. In contrast, OFT employs a layer-wise, shared orthogonal transformation, multiplicatively modifying neuron weights. This formulation mathematically guarantees the preservation of pairwise angles among neurons within each layer, leading to enhanced training stability.

\section{Orthogonal Finetuning}

\subsection{Why Does Orthogonal Transformation Make Sense?}

To motivate the use of orthogonal transformation in fine-tuning text-to-image diffusion models, we break the question into two parts: (1) the rationale for focusing on modifying neuron angles (i.e., directions) during fine-tuning, and (2) the justification for using orthogonal transformations to effect these angle changes.

Addressing the first question, our approach draws from prior empirical findings in \cite{liu2018decoupled,chen2020angular}, which indicate that angular differences among features are strongly indicative of the semantic gap. SphereNet~\cite{liu2017deep} demonstrates that confining neural activations to lie on a unit hypersphere preserves model expressivity and often leads to improved generalization, suggesting that the directional component of neurons encapsulates the core data-relevant information. To further clarify the significance of angular neuron representations, we present a simple illustrative experiment in Figure~\ref{fig:toy_ae}, where a standard convolutional autoencoder is trained using flower image data. During training, feature maps are computed via the conventional inner product (where $z$ is the convolution kernel $\bm{w}$ applied to input patch $\bm{x}$). At test time, we consider three approaches for obtaining feature maps: (a) reusing the inner product as in training, (b) extracting only the magnitude, and (c) using purely the angle. The reconstructions displayed in Figure~\ref{fig:toy_ae} reveal that the angular component alone is almost sufficient to reconstruct the original images, whereas the magnitudes carry little to no relevant information. Importantly, cosine-based activations (case (c)) are not employed in training—only the inner product is used—underscoring that semantic information is fundamentally captured in neuron orientations. Consequently, to alter the semantic content of images, adjusting neuron directions (angles) is likely more effective than manipulating magnitudes.

For the second question, a straightforward method to refine neuron directions involves applying a uniform rotation or reflection to all neurons within the same layer, which is naturally represented by an orthogonal transformation. Although more complex angular modifications—such as rotating individual neurons by distinct angles—are possible, orthogonal transformations offer a favorable balance between adaptability and structure. Furthermore, results from \cite{liu2021orthogonal} indicate that orthogonal transformations possess ample expressive capacity for neural network learning. To illustrate the practical benefits of orthogonal constraints, we replicate a controlled generation task following the ControlNet~\cite{zhang2023adding} protocol, with outcomes shown in Figure~\ref{fig:ortho}. Our method, employing standard OFT, yields stable training and precise control by the conclusion of training (epoch 20). In contrast, removing the orthogonality constraint causes the model to fail in generating realistic or controllable images, emphasizing the crucial regularization effect provided by orthogonality in OFT.

\subsection{General Framework}

Traditional finetuning generally relies on gradient descent with a modest learning rate to modify the parameters of a model, or specific layers within it. The use of a low learning rate effectively limits the magnitude of changes made to the pretrained weights, which can be viewed as a form of regularization. Specifically, classical finetuning attempts to learn model weights $\bm{M}$ while indirectly restricting the distance from the initial, pretrained weights $\bm{M}^0$, as measured by $\thickmuskip=2mu \medmuskip=2mu \| \bm{M} - \bm{M}^0\|$. While this inherent regularization is intuitively appealing, it may still allow for excessive flexibility, particularly when finetuning very large models. To mitigate this, LoRA applies an explicit low-rank limitation on weight updates, specified as $\thickmuskip=2mu \medmuskip=2mu \text{rank}(\bm{M} - \bm{M}^0)=r'$, where $r'$ is chosen to be small. In contrast, OFT introduces a different form of constraint by enforcing pairwise neuron similarity: $\thickmuskip=2mu \medmuskip=2mu \| \text{HE}(\bm{M})-\text{HE}(\bm{M}^0) \|=0$, with $\text{HE}(\cdot)$ representing the hyperspherical energy of the weight matrix. 

To illustrate this concept, suppose we examine a fully connected layer defined by the weight matrix $\thickmuskip=2mu \medmuskip=2mu \bm{W}=\{\bm{w}_1,\cdots,\bm{w}_n\}\in\mathbb{R}^{d\times n}$, where each $\bm{w}_i\in\mathbb{R}^{d}$ corresponds to the $i$-th neuron, and $\bm{W}^0$ is its pretrained counterpart. The output $\thickmuskip=2mu \medmuskip=2mu \bm{z}\in\mathbb{R}^n$ produced by this layer given an input $\thickmuskip=2mu \medmuskip=2mu \bm{x}\in\mathbb{R}^d$ is calculated as $\thickmuskip=2mu \medmuskip=2mu \bm{z}=\bm{W}^\top\bm{x}$. From this perspective, OFT may be interpreted as striving to minimize the difference in hyperspherical energy between the finetuned and pretrained weights:

\begin{equation}\label{eq:oft_principle}
\footnotesize
\min_{\bm{W}} \left\| \text{HE}(\bm{W})-\text{HE}(\bm{W}^0)\right\|~~\Leftrightarrow~~\min_{\bm{W}} \bigg{\|} \sum_{i\neq j} \|\hat{\bm{w}}_i-\hat{\bm{w}}_j\|^{-1}-\sum_{i\neq j} \|\hat{\bm{w}}^0_i-\hat{\bm{w}}^0_j\|^{-1}\bigg{\|}
\end{equation}

Here, $\thickmuskip=2mu \medmuskip=2mu \hat{\bm{w}}_i:=\bm{w}_i/\|\bm{w}_i\|$ denotes the normalized version of neuron $i$, and the hyperspherical energy for a layer $\bm{W}$ is computed as $\thickmuskip=2mu \medmuskip=2mu \text{HE}(\bm{W}):= \sum_{i\neq j} \|\hat{\bm{w}}_i-\hat{\bm{w}}_j\|^{-1}$. It is straightforward to verify that the global minimum of Eq.~\eqref{eq:oft_principle} is zero, and this minimum is reached whenever $\bm{W}$ is related to $\bm{W}^0$ by an orthogonal transformation; that is, $\thickmuskip=2mu \medmuskip=2mu \bm{W}=\bm{R}\bm{W}^0$ where $\thickmuskip=2mu \medmuskip=2mu \bm{R}\in\mathbb{R}^{d\times d}$ is an orthogonal matrix. Depending on the determinant of $\bm{R}$ (either $1$ or $-1$),

\begin{equation}\label{eq:oft_general}
    \footnotesize
    \bm{z}=\bm{W}^\top\bm{x}=(\bm{R}\cdot\bm{W}^0)^\top\bm{x},~~~\text{s.t.}~\bm{R}^\top\bm{R}=\bm{R}\bm{R}^\top=\bm{I}
\end{equation}

Here, $\bm{W}$ refers to the weight matrix after OFT-based finetuning, and $\bm{I}$ stands for the identity matrix. Figure~\ref{fig:block} provides a depiction of the OFT framework. In analogy to the approach used by LoRA with zero initialization, we must guarantee that OFT performs finetuning starting strictly from the initial pretrained parameters $\bm{W}^0$. To facilitate this, the orthogonal matrix $\bm{R}$ is initially set as the identity, thereby ensuring the initial state matches the pretrained model. Maintaining the orthogonality property of $\bm{R}$ is critical; for this purpose, we can employ differentiable orthogonalization techniques as described in \cite{liu2021orthogonal,lezcano2019cheap}. We elaborate later on efficient computational strategies for preserving orthogonality.

\subsection{Efficient Orthogonal Parameterization}\label{sect:eff_ortho}

Although classic orthogonalization techniques, such as the Gram-Schmidt procedure, are differentiable, their computational demands make them impractical for large-scale applications~\cite{liu2021orthogonal}. To enhance efficiency, we utilize the Cayley transform to parameterize the orthogonal matrix. More precisely, we define $\bm{R}=(\bm{I}+\bm{Q})(I-\bm{Q})^{-1}$, where $\bm{Q}$ is a skew-symmetric matrix, satisfying $\bm{Q}=-\bm{Q}^\top$. The Cayley transform accelerates computation but restricts us to generating orthogonal matrices with determinant $1$, that is, matrices from the special orthogonal group. Empirically, we observe this restriction does not impair performance. Despite using the Cayley parameterization, representing the entire orthogonal matrix $\bm{R}$ may still be inefficient when the dimensionality $d$ is large. To overcome this, we suggest structuring $\bm{R}$ as a block-diagonal matrix comprising $r$ blocks, leading to the following structure:

\begin{equation}
    \footnotesize
    \bm{R}=\text{diag}(\bm{R}_1,\bm{R}_2,\cdots,\bm{R}_r)=\begin{bmatrix}
    \bm{R}_1\in \text{O}(\frac{d}{r}) &  &\\
    &\ddots & \\
    & & \bm{R}_r\in \text{O}(\frac{d}{r})
    \end{bmatrix}\in \text{O}(d)
\end{equation}

Here, $\text{O}(d)$ represents the orthogonal group of $d$ dimensions. The matrix $\thickmuskip=2mu \medmuskip=2mu \bm{R}\in\mathbb{R}^{d\times d}$ and each block $\thickmuskip=2mu \medmuskip=2mu \bm{R}_i\in\mathbb{R}^{d/r\times d/r}$ for all $i$ are orthogonal. When we set $\thickmuskip=2mu \medmuskip=2mu r=1$, the block-diagonal construction simply reduces to a general unconstrained orthogonal matrix. For a full orthogonal matrix of dimensions $\thickmuskip=2mu \medmuskip=2mu d\times d$, the number of independent parameters is $\thickmuskip=2mu \medmuskip=2mu d(d-1)/2$, leading to computational complexity $\mathcal{O}(d^2)$. In contrast, an orthogonal matrix composed of $r$ blocks requires only $\thickmuskip=2mu \medmuskip=2mu d(d/r-1)/2$ parameters, yielding complexity $\thickmuskip=2mu \medmuskip=2mu \mathcal{O}(d^2/r)$. To minimize parameters even more, one can opt for parameter sharing among blocks, \emph{i.e.}, enforcing $\thickmuskip=2mu \medmuskip=2mu\bm{R}_i=\bm{R}_j$ for all $i\neq j$. This reduction changes the parameter count to $\mathcal{O}(d^2/r^2)$. Importantly, despite these parameter-saving schemes, the composed matrix $\bm{R}$ retains its orthogonality, thereby fully maintaining the hyperspherical structure.

To compare the parameter efficiency of OFT against that of LoRA, consider LoRA’s formulation with a low-rank parameter $r'$, which produces $r'(d+n)$ trainable parameters. If we suppose $r$ and $r'$ both scale with the hidden size $d$ (\emph{e.g.}, $r = r' = \alpha d$, with some fixed $0 < \alpha \leq 1$), then LoRA’s parameter complexity is $\mathcal{O}(d^2 + dn)$, while OFT achieves $\mathcal{O}(d)$ complexity. A practical example clarifies this difference: for a weight matrix sized $128 \times 128$, LoRA with $r' = 8$ results in $2,048$ parameters, whereas OFT with $r=8$ (with no sharing across blocks) uses only $960$ parameters.

\subsection{Constrained Orthogonal Finetuning}

To further tighten OFT’s parameter flexibility, we propose restricting the finetuned model to remain close to the initial pretrained model. COFT implements this idea via the following forward computation:

\begin{equation}\label{eq:coft_general}
    \footnotesize
    \bm{z}=\bm{W}^\top\bm{x}=(\bm{R}\cdot\bm{W}^0)^\top\bm{x},~~~\text{s.t.}~\bm{R}^\top\bm{R}=\bm{R}\bm{R}^\top=\bm{I},~\norm{\bm{R}-\bm{I}}\leq \epsilon
\end{equation}

where an orthogonality constraint and an additional proximity constraint ($\epsilon$-bound from the identity) are enforced. The orthogonality property can be realized through Cayley parameterization, as explained in Section~\ref{sect:eff_ortho}. However, integrating the $\epsilon$-deviation directly with a Cayley-parameterized matrix is challenging. To investigate this further, we apply the Neumann expansion to approximate $\thickmuskip=2mu \medmuskip=2mu \bm{R}=(\bm{I}+\bm{Q})(I-\bm{Q})^{-1}$ as $\thickmuskip=2mu \medmuskip=2mu \bm{R}\approx\bm{I}+2\bm{Q}+\mathcal{O}(\bm{Q}^2)$, assuming that the Neumann

Given that the series converges with respect to the operator norm, it becomes possible to shift the restriction $\thickmuskip=2mu \medmuskip=2mu\|\bm{R}-\bm{I}\|\leq\epsilon$ inside the Cayley transform. The associated equivalent constraint is $\thickmuskip=2mu \medmuskip=2mu \|\bm{Q}-\bm{0}\|\leq\epsilon'$, where $\bm{0}$ signifies the zero matrix and $\epsilon'$ serves as a different tolerance parameter from $\epsilon$. This reformulated restriction on $\bm{Q}$ is readily enforced using projected gradient descent techniques. To obtain identity initialization for the orthogonal transformation $\bm{R}$, one initializes $\bm{Q}$ to be a zero matrix. Conceptually, COFT imposes two explicit constraints: bounded deviation in hyperspherical energy and an explicit bound on the departure from the pretrained model parameters. While many finetuning approaches typically rely on the latter as an implicit condition, COFT highlights and enforces it directly. Although OFT already provides strong results, we note that the explicit deviation constraint in COFT enhances finetuning stability even further. Figure~\ref{fig:eps_coft} illustrates how the hyperparameter $\epsilon$ modulates COFT’s behavior; specifically, increasing $\epsilon$ allows the COFT-adapted model to more closely approximate the model fine-tuned by OFT, whereas decreasing $\epsilon$ causes the result to resemble the original pretrained text-to-image diffusion model.

\subsection{Re-scaled Orthogonal Finetuning}

We introduce a straightforward extension to OFT whereby each neuron acquires its own learnable scaling coefficient. The rationale behind this lies in the invariance of hyperspherical energy to such scaling, as normalization offsets the effect of magnitude changes. Concretely, the forward computation is modified to: $\thickmuskip=2mu \medmuskip=2mu\bm{z}=(\bm{R}\bm{W}^0\bm{D})^\top\bm{x}$\footnote{\textbf{Errata}: In the NeurIPS camera ready paper, the re-scaled OFT forward computation was incorrectly typeset as $\thickmuskip=2mu \medmuskip=2mu\bm{z}=(\bm{D}\bm{R}\bm{W}^0)^\top\bm{x}$. The implementation used is correct; results in Appendix~\ref{app:rescaled} remain valid.}, where $\thickmuskip=2mu \medmuskip=2mu\bm{D}=\text{diag}(s_1,\cdots,s_n)\in\mathbb{R}^{n\times n}$ represents a diagonal matrix whose entries $s_1,\ldots,s_n$ are learnable and strictly positive. Unlike the original OFT’s approach, where only $\bm{R}$ is trainable as in Eq.~\eqref{eq:oft_general}, both $\bm{D}$ and $\bm{R}$ are updated during training in this variant. This re-scaled OFT achieves additional adaptability at a marginal computational cost. To isolate the effectiveness of orthogonal transformation, our main experiments use unmodified OFT; nevertheless, in practice, we observe the re-scaled variant often yields superior performance (as evidenced in Appendix~\ref{app:rescaled}).

\section{Intriguing Insights and Discussions}

\textbf{OFT applies across diverse architectures}. OFT is, in theory, applicable to any neural network model. Within Transformer architectures, LoRA is customarily utilized for updating the attention weight matrices~\cite{hulora2022}. For direct comparison with LoRA, we likewise confine OFT to only finetuning the attention weights throughout our study. Beyond its utility for fully connected layers, OFT is also readily extendable to convolutional layers, due to the interpretable structure that the block-diagonal $\bm{R}$ matrix imparts in this context, setting it apart from LoRA. Specifically, assigning the number of blocks in $\bm{R}$ to match the input channel count allows each block to affect a distinct neuron channel, analogous to depth-wise convolution kernel adaptation~\cite{chollet2017xception}. If $\bm{R}$'s blocks are shared, all neurons experience a common orthogonal transformation across channels. Appendix~\ref{app:convolution} reports early results exploring OFT for convolutional layer finetuning.

\textbf{Connection to LoRA}. LoRA introduces a low-rank update to the pretrained weight matrix, thereby constraining drastic alterations to its underlying information. In comparison, OFT maintains the integrity of pretraining by enforcing the transformation applied to the weight matrix to be orthogonal (i.e., full-rank), which guards against the loss of pretrained knowledge. The forward propagation in OFT can equivalently be expressed as $\thickmuskip=2mu \medmuskip=2mu \bm{z}=(\bm{R}\bm{W}^0)^\top\bm{x}=(\bm{W}^0+(\bm{R}-\bm{I})\bm{W}^0)^\top\bm{x}$, where $\thickmuskip=2mu \medmuskip=2mu (\bm{R}-\bm{I})\bm{W}^0$ serves as a counterpart to LoRA’s low-rank modification. Since $\bm{W}^0$ is generally of full rank, OFT, too, executes a low-rank adaptation whenever $\thickmuskip=2mu \medmuskip=2mu \bm{R}-\bm{I}$ exhibits low rank. Similar to how LoRA uses a rank hyperparameter $r'$ to manage trainable parameter counts, OFT relies on a block-diagonal parameter $r$ for parameter reduction. Notably, while both techniques increase parameter efficiency, they pursue distinct strategies: LoRA reduces the dimensionality of updates through low-rank structure, whereas OFT leverages block-diagonal orthogonality, introducing sparsity into the trainable parameter space.

\textbf{Why OFT converges faster}. The evidence from Figure~\ref{fig:toy_ae} indicates that semantic modification is most effectively achieved via adjustments in neuron direction, a property that OFT is explicitly engineered to exploit. Moreover, OFT’s optimization process can be interpreted as adapting neurons on the smooth surface of a hypersphere, yielding a more favorable optimization landscape. This observation aligns with findings reported in \cite{liu2021orthogonal}.

\textbf{Why not minimize hyperspherical energy}. In contrast with \cite{liu2021orthogonal}, our approach does not involve minimizing the hyperspherical energy. In classification settings, it is beneficial for neurons to be non-redundant; arranging neurons to minimize hyperspherical energy would space them evenly across the hypersphere. However, this uniform arrangement could eliminate important information acquired during pretraining, making such minimization unsuitable for the finetuning context.

\textbf{Balancing flexibility and regularity in finetuning}. Our findings indicate that there is an inherent balance to be struck between flexibility and regularity in finetuning approaches. While conventional finetuning techniques offer high flexibility, they tend to suffer from instability and are prone to issues like model collapse. The OFT method, though straightforward, effectively mediates this trade-off by maintaining pairwise neuron angles, striking a favorable compromise between adaptability and structure. Additionally, utilizing a block-diagonal approach to parameterize the orthogonal matrix serves to further regularize the matrix in a more stringent manner.

\textbf{No increase in inference computation}. In contrast to methods like ControlNet, our OFT framework does not add any extra computational load during inference for the finetuned model. Specifically, during the inference phase, one can directly incorporate the trained orthogonal matrix $\bm{R}$ into the original pretrained weights $\bm{W}^0$ through simple multiplication, yielding $\thickmuskip=2mu \medmuskip=2mu \bm{W}=\bm{R}\bm{W}^0$, an equivalent parameter matrix. This ensures that the inference performance, in terms of speed, remains identical to that of the original pretrained network.

\section{Experiments and Results}

\textbf{Overall experimental settings}. For our experimental setup, we employ Stable Diffusion v1.5~\cite{rombach2022high} as the backbone text-to-image generation model. To ensure fair comparison, image samples for each method are selected at random. For subject-driven generation tasks, the experimental protocol follows DreamBooth~\cite{ruiz2023dreambooth}, while controllable generation experiments align with the methodologies in ControlNet~\cite{zhang2023adding} and T2I-Adapter~\cite{mou2023t2i}. In order to fairly benchmark against LoRA, the application of OFT or COFT is restricted to the same layers where LoRA is utilized. Additional experimental details and supplementary results can be found in Appendix~\ref{app:settings}.

\subsection{Subject-driven Generation}

\textbf{Experimental details}. Our baseline methods comprise DreamBooth~\cite{ruiz2023dreambooth} and LoRA~\cite{hulora2022}. All compared approaches utilize the same loss as defined in DreamBooth. For the baselines, we adhere to the protocols and optimal hyperparameters described in their respective works. Further results and extensive experimental details are available in Appendix~\ref{app:settings},~\ref{app:coft_oft},~\ref{app:more_qualitative}, and~\ref{app:failure}.

\textbf{Finetuning stability and convergence}. We begin by assessing both the stability of finetuning and the convergence rates of DreamBooth, LoRA, OFT, and COFT. These comparisons are visualized in Figure~\ref{fig:teaser} and Figure~\ref{fig:db_convergence}. Our findings indicate that the diffusion models trained using either COFT or OFT demonstrate highly stable finetuning dynamics. By 400 training steps, DreamBooth and both OFT-based approaches manage to effectively steer the generative process, whereas LoRA does not maintain accurate subject identity. Extending training to 2000 iterations leads DreamBooth to collapse—producing degenerate outputs—and LoRA to fail in generating the correct clothing color (producing “yellow fur” instead of “yellow shirt”). In stark contrast, both OFT and COFT continue to provide robust and reliable control of the output images. These observations affirm that our OFT framework supports rapid, yet stable convergence for subject-driven generation tasks. The property of being largely insensitive to the number of finetuning steps is especially valuable, as it eliminates the need for subject-specific tuning of iteration counts. Both OFT and COFT permit the use of a larger iteration number as a default, circumventing exhaustive hyperparameter search. Notably, for COFT—when paired with a suitable $\epsilon$—neither learning rate nor the number of training iterations require fine adjustment.

\textbf{Quantitative comparison}. Adopting the evaluation framework from \cite{ruiz2023dreambooth}, we perform quantitative assessments targeting subject fidelity (using DINO~\cite{caron2021emerging} and CLIP-I~\cite{radford2021learning}), text prompt alignment (using CLIP-T~\cite{radford2021learning}), and generative diversity (using LPIPS~\cite{zhang2018unreasonable}). The CLIP-I metric measures the average cosine similarity between CLIP embeddings derived from the generated samples and actual images. The DINO metric follows the same logic, but leverages ViT S/16 DINO embeddings. The CLIP-T score reflects the mean cosine similarity between embeddings of the prompt text and the corresponding generated images. To gauge diversity, we also compute the average LPIPS cosine similarity among images synthesized for the same prompt and subject. The empirical results, shown in Table~\ref{table:dreambooth_final}, demonstrate that both COFT and OFT achieve notably higher DINO and CLIP-I scores than DreamBooth and LoRA, with competitive or slightly improved outcomes in terms of prompt fidelity and diversity. For every method, we repeat finetuning from the same pretrained model across 30 distinct random seeds to reduce the influence of stochastic variation. Collectively, the results indicate that the OFT framework not only excels in terms of stability and convergence, but also consistently delivers superior final performance.

\textbf{Qualitative comparison}. To provide a clearer picture of OFT’s strengths, Figure~\ref{fig:dreambooth_final} presents a selection of randomly chosen samples illustrating subject-driven generation. All examples derive from the same finetuned base model for each approach, ensuring consistency as no specific prompt optimizations are applied to individual outputs. For every method, we utilize the model checkpoint yielding the optimal validation CLIP metrics. According to the observations in Figure~\ref{fig:dreambooth_final}, OFT and COFT both excel in faithfully preserving the semantic features of the target subject, in contrast to LoRA, which frequently struggles to maintain subject identity (\emph{e.g.}, LoRA completely fails to retain the identity of the bowl in the corresponding case). Furthermore, OFT and COFT exhibit greater precision when conditioned on textual prompts, whereas DreamBooth, although generally successful at retaining subject identity, is less reliable in accurately interpreting prompt instructions (\emph{e.g.}, the top-row bowl instance). These results indicate that the OFT framework effectively combines both prompt controllability and robust subject preservation. Additionally, OFT’s performance remains stable across a range of iteration counts, enabling robust results even at high iteration numbers—a property not shared by either DreamBooth or LoRA. Additional examples can be found in Appendix~\ref{app:more_qualitative}. In addition, a human assessment conducted in Appendix~\ref{app:human} offers further support for the advantages of OFT.

\subsection{Controllable Generation}

\textbf{Settings}. We adopt ControlNet~\cite{zhang2023adding}, T2I-Adapter~\cite{mou2023t2i}, and LoRA~\cite{hulora2022} as comparison baselines. The main paper investigates three challenging controllable generation scenarios: canny edge to image~(C2I) with the COCO dataset~\cite{lin2014microsoft}; segmentation map to image~(S2I) on ADE20K~\cite{zhou2017scene}; and landmark to face~(L2F) with CelebA-HQ~\cite{karras2018progressive,xia2021tedigan}. For each task, all methods finetune Stable Diffusion~(SD) v1.5 for 20 training epochs. Expanded experimental results appear in Appendix~\ref{app:more_qualitative},\ref{app:more_control_task},\ref{app:failure}.

\textbf{Convergence}. To assess the rate of convergence, we analyze ControlNet, T2I-Adapter, LoRA, and COFT on the L2F task, employing both quantitative metrics and qualitative assessments. Quantitatively, we report the mean $\ell_2$ distance between the target control face landmarks and the predicted face landmarks as our primary evaluation metric. Figure~\ref{fig:face_convergence} presents the error in face landmarks as a function of the number of fine-tuning epochs for each model. The results demonstrate that COFT and OFT converge at a substantially higher rate than the baselines; for example, LoRA requires 20 epochs to attain the same level of performance that OFT achieves at epoch 8. Importantly, both OFT and COFT have a comparable count of trainable parameters to LoRA (significantly less than ControlNet), yet they display much more rapid convergence relative to existing alternatives. Additional evidence of OFT's swift convergence can be found in Figure~\ref{fig:teaser}, where the rightmost example illustrates that OFT achieves effective convergence using just 5\% of the ADE20K dataset—highlighting a marked improvement in data efficiency over ControlNet and LoRA. Qualitative analysis centers on OFT, COFT, and ControlNet due to ControlNet exhibiting a comparable landmark error. As seen in Figure~\ref{fig:face_convergence_qual}, both OFT and COFT not only converge steadily but also show consistent improvement in the alignment of generated face poses to the reference landmarks across epochs. By contrast, ControlNet's controllability exhibits an abrupt improvement post-8 epochs, reflecting the "sudden convergence" previously described in \cite{zhang2023adding}. This smooth and stable convergence trajectory makes the OFT framework highly practical, as its training dynamics are more reliable and straightforward to interpret, ultimately simplifying the process of tuning OFT’s hyperparameters.

\textbf{Quantitative comparison}. To assess the efficacy of controllable generation, we propose a metric that quantifies control consistency. This involves extracting the control signal from the synthesized image and directly comparing it with the original input signal. Specifically, for the C2I task, we report IoU and F1 scores; for the S2I task, we calculate mean IoU as well as both mean and overall accuracy; and for the L2F task, we determine the average $\ell_2$ distance between the predicted and target control landmarks. Additional specifics on these consistency metrics are provided in Appendix~\ref{app:settings}. All methods under comparison were evaluated using their optimal hyperparameter configurations. According to Table~\ref{table:control_gen}, both OFT and COFT consistently achieve stronger and more precise control compared to alternative methods. Notably, adapter-based solutions (such as T2I-Adapter and ControlNet) exhibit slower convergence rates and inferior ultimate performance. While LoRA demonstrates superior outcomes relative to ControlNet on the S2I task, it underperforms on both C2I and L2F tasks. Overall, the maximum performance attainable with LoRA appears constrained, even after extensive hyperparameter tuning. In contrast, OFT continues to improve with increasing trainable parameter counts, suggesting its performance has yet to plateau. Importantly, our systematic quantitative assessment of controllable generation constitutes a key contribution of this work, providing an accurate measure of control efficacy in finetuned models—subject to the inherent accuracy of the underlying segmentation or detection modules.

\textbf{Qualitative comparison}. We further conduct a qualitative assessment of OFT and COFT alongside leading contemporary methods such as ControlNet, T2I-Adapter, and LoRA. As illustrated by the randomly generated samples in Figure~\ref{fig:control_final}, both OFT and COFT are capable of producing images that exhibit not only photorealistic quality but also precise adherence to the intended control inputs. For the S2I task, it is evident that LoRA is unable to accurately reflect the input segmentation maps in its outputs, whereas ControlNet, OFT, and COFT effectively respect the segmentation guidance. Notably, OFT and COFT surpass ControlNet by producing images with enhanced detail and more plausible geometry, despite utilizing considerably fewer model parameters. Regarding the C2I task, OFT and COFT successfully generate semantically coherent images from rough Canny edge inputs, outperforming T2I-Adapter and LoRA, which demonstrate inferior control. In the L2F task, our approach achieves the most precise correspondence between input poses and generated faces, even under complex pose conditions. Across all three control tasks, OFT and COFT consistently deliver images of superior qualitative quality compared to state-of-the-art baselines, attesting to the efficacy of our OFT architecture for controllable generation. For broader qualitative analysis, we present additional results spanning the three control tasks in Appendix~\ref{app:control_exp}, and we further show the versatility of OFT on additional control challenges (such as dense pose to human body, sketch to image, and depth to image) in Appendix~\ref{app:more_control_task}.

\section{Concluding Remarks and Open Problems}

Prompted by the critical role angular relationships among neurons play in shaping visual semantics, we have introduced a straightforward yet powerful approach—orthogonal finetuning—for text-to-image diffusion model control. In particular, our method is directed at two primary tasks: subject-driven generation and controllable generation. Relative to previous techniques, OFT offers enhanced controllability and more stable finetuning performance, all while requiring fewer parameters for adaptation. Furthermore, OFT preserves inference efficiency by not adding any additional computational burden, making it practical for deployment.

Nevertheless, OFT opens the door to several compelling unresolved questions. Firstly, the method enforces orthogonality through Cayley parametrization, which necessitates a matrix inversion—a process that somewhat constrains the scalability of OFT. While adopting a block diagonal parametrization helps alleviate this, devising ways to accelerate this inversion operation in a differentiable context remains a significant hurdle. Secondly, OFT presents promising opportunities in the realm of compositionality; specifically, the orthogonal matrices generated through separate OFT finetuning sessions can be composed via multiplication, yielding another orthogonal matrix. It remains to be thoroughly investigated whether this ensemble of orthogonal matrices can simultaneously retain expertise from all the corresponding finetuning tasks. Finally, OFT's parameter efficiency currently relies heavily on its block diagonal structure, which, unfortunately, may introduce certain biases and constrain flexibility. Exploring strategies to further boost parameter efficiency without incurring extra bias or sacrificing adaptability stands as a notable challenge for future research.

\end{document}